\documentclass[../DoAn.tex]{subfiles}

\begin{document}
\sloppy

Trước khi thiết kế bất cứ thứ gì, tôi cần trả lời một câu hỏi tưởng đơn giản:
vấn đề thực sự của người dùng là gì? Câu hỏi đó dẫn đến nhiều cuộc trao đổi
với sinh viên đang ở ký túc xá, sinh viên đang chờ kết quả phân bổ, và cán bộ
đang xử lý hàng trăm hồ sơ mỗi đợt. Kết quả thu được không hoàn toàn như tôi
kỳ vọng ban đầu --- vấn đề không chỉ là quy trình chậm mà còn là sự thiếu
tin tưởng: sinh viên không biết mình được đánh giá theo tiêu chí gì, cán bộ
không có công cụ để kiểm chứng quyết định trước khi nó có hiệu lực.

Dựa trên những phát hiện đó, chương này khảo sát thực trạng vận hành, phân
tích nhu cầu thực tế của các đối tượng sử dụng, và nghiên cứu ưu nhược điểm
của các hệ thống quản lý ký túc xá hiện có, từ đó xác định khoảng trống mà
giải pháp đề xuất hướng tới lấp đầy. Nội dung được tổ chức thành bốn phần:
khảo sát hiện trạng và xu hướng phát triển (2.1), tổng quan chức năng và biểu
đồ use case (2.2), đặc tả các trường hợp sử dụng trọng yếu (2.3) và yêu cầu
phi chức năng (2.4).

% =====================================================================
\section{Khảo sát hiện trạng}
\label{sec:survey}
% =====================================================================

\subsection{Khảo sát người dùng}

Để xác định các yêu cầu thực tế của hệ thống, tôi đã trao đổi với sinh viên đang
cư trú tại ký túc xá, sinh viên tham gia đăng ký chỗ ở và cán bộ quản lý. Hoạt
động khảo sát được thực hiện dưới hình thức phỏng vấn mở nhằm thu thập những
khó khăn và nhu cầu phát sinh trong quá trình sử dụng hệ thống hiện tại.

Kết quả khảo sát cho thấy nhiều sinh viên gặp khó khăn trong việc tiếp cận thông
tin liên quan đến quá trình phân bổ chỗ ở. Người dùng thường không hiểu rõ tiêu
chí xét duyệt, cơ chế ưu tiên và nguyên nhân dẫn đến kết quả phân bổ. Bên cạnh
đó, sinh viên cũng không thể hình dung rõ môi trường lưu trú trước khi đăng ký
do thiếu các thông tin trực quan về vị trí phòng, sơ đồ tầng và cơ sở vật chất
thực tế. Điều này khiến quá trình lựa chọn chỗ ở phụ thuộc nhiều vào kinh
nghiệm truyền miệng từ các khóa trước.

Các ý kiến thu thập được có thể được phân thành ba nhóm vấn đề chính. Thứ nhất,
quy trình phân bổ còn thiếu tính minh bạch, khiến sinh viên khó hiểu và khó
kiểm chứng kết quả xét duyệt. Thứ hai, trải nghiệm đăng ký còn hạn chế do thiếu
cơ chế theo dõi trạng thái hồ sơ và phản hồi theo thời gian thực. Thứ ba, cán
bộ quản lý chưa có công cụ hỗ trợ mô phỏng chính sách và đánh giá tác động trước
khi triển khai trên dữ liệu thực tế.

Ngoài ra, nhu cầu xem trước không gian phòng ở được ghi nhận với tần suất cao.
Nhiều sinh viên mong muốn được tham khảo hình ảnh thực tế, sơ đồ tầng hoặc mô
hình trực quan trước khi đưa ra quyết định đăng ký. Người dùng cũng kỳ vọng
nhận được thông báo kịp thời khi hồ sơ thay đổi trạng thái hoặc khi kết quả
phân bổ được công bố. Từ những nhu cầu này, trực quan hóa phòng ở, phân bổ tự
động, mô phỏng chính sách và thông báo thời gian thực được xác định là các yêu
cầu cốt lõi của hệ thống.

Trong các vấn đề được ghi nhận, tính minh bạch của quy trình phân bổ được đánh
giá là yếu tố quan trọng nhất. Khi sinh viên không hiểu hoặc không tin tưởng
vào kết quả phân bổ, các cải tiến về giao diện hay hiệu năng hệ thống sẽ khó
mang lại hiệu quả thực sự. Vì vậy, khả năng giải thích kết quả và truy vết quá
trình phân bổ được xem là một trong những nguyên tắc thiết kế trọng tâm của dự
án.


\subsection{Xu hướng phát triển các hệ thống quản lý ký túc xá}

Trước khi phân tích xu hướng thị trường, cần nhìn nhận thực trạng của phần lớn
hệ thống quản lý ký túc xá hiện nay tại Việt Nam. Thông tin về phòng ở thường
bị phân tán qua nhiều kênh --- bảng thông báo, email, website tĩnh --- khiến
sinh viên khó tổng hợp để đưa ra quyết định. Quy trình đăng ký vẫn phụ thuộc
nhiều vào thủ tục giấy tờ hoặc biểu mẫu điện tử đơn giản, thiếu khả năng theo
dõi trạng thái hồ sơ theo thời gian thực. Đặc biệt, sinh viên không có cơ hội
xem trước không gian phòng trước khi đăng ký, dẫn đến quyết định thiếu căn cứ
và gia tăng số lượng điều chỉnh sau phân bổ.

Một hạn chế khác thường bị bỏ qua là khả năng tiếp cận đối với sinh viên quốc
tế. Các hệ thống đang vận hành chủ yếu chỉ hỗ trợ tiếng Việt và giả định người
dùng quen với quy trình hành chính trong nước. Điều này tạo ra rào cản không
nhỏ cho sinh viên trao đổi, du học sinh và sinh viên quốc tế --- nhóm đối tượng
ngày càng chiếm tỷ trọng lớn hơn trong cộng đồng đại học, đặc biệt tại các
trường có định hướng hội nhập quốc tế như Đại học Bách Khoa Hà Nội.

Để hiểu rõ hơn xu hướng phát triển của thị trường và tìm kiếm các hướng tiếp
cận tham khảo, tôi tiến hành khảo sát một số nền tảng quản lý chỗ ở sinh viên
quốc tế đang hoạt động hiệu quả, từ đó rút ra bài học cho việc thiết kế hệ
thống eDorm.

\subsubsection*{Khảo sát nền tảng Scape}

Scape (\textit{scape.com.au}) là nền tảng cho thuê chỗ ở sinh viên tại Úc, nổi
bật với định hướng lấy người dùng làm trung tâm trong toàn bộ hành trình tìm
kiếm và đăng ký. Điểm đáng chú ý nhất là tính năng tham quan phòng bằng tour
3D 360° tích hợp ngay trong trang chi tiết, cho phép sinh viên trải nghiệm
không gian thực tế trước khi đưa ra quyết định mà không cần đến trực tiếp. Hệ
thống hiển thị tình trạng phòng còn trống theo thời gian thực và hỗ trợ lựa
chọn thời gian nhận phòng linh hoạt theo học kỳ. Thông tin về loại phòng, số
người ở, tiện nghi đi kèm, giá thuê và bộ lọc giới tính được trình bày rõ ràng
theo từng danh mục, giúp sinh viên tự tìm hiểu và so sánh mà không cần liên hệ
trực tiếp với nhân viên.

\begin{figure}[H]
\centering
\includegraphics[width=0.6\linewidth]{Hinhve/khaosat_scape.png}
\caption{Giao diện khảo sát từ nền tảng Scape.com.au}
\label{fig:scape}
\end{figure}

Từ góc độ trải nghiệm người dùng, Scape thể hiện rõ mô hình tự phục vụ khi sinh
viên có thể hoàn thành toàn bộ quy trình từ tìm kiếm, so sánh đến đăng ký chỗ ở
mà không cần tương tác trực tiếp với nhân viên hỗ trợ. Giao diện được tối ưu
cho cả máy tính và thiết bị di động, tốc độ phản hồi nhanh và số lượng thao tác
được rút gọn tối đa. Đây là điểm đáng tham khảo đối với các hệ thống quản lý ký
túc xá trong nước, nơi quy trình đăng ký vẫn còn phụ thuộc vào nhiều bước xử lý
và xác nhận thủ công.

Tuy nhiên, mô hình vận hành của Scape có sự khác biệt cơ bản so với bài toán
quản lý ký túc xá tại các trường đại học công lập ở Việt Nam. Trong khi người
dùng của Scape chủ động lựa chọn phòng dựa trên nhu cầu cá nhân và khả năng tài
chính, sinh viên đăng ký ký túc xá thường được phân bổ chỗ ở theo các tiêu chí
và chính sách ưu tiên do nhà trường quy định. Vì vậy, quy trình của Scape mang
tính lựa chọn dịch vụ, trong khi quy trình ký túc xá mang tính xét duyệt và phân
bổ nguồn lực.

Do đó, giá trị tham khảo chính từ Scape không nằm ở cơ chế ra quyết định mà ở
cách trình bày thông tin và trực quan hóa không gian phòng ở. Đặc biệt, chức
năng tham quan phòng trực tuyến đã cung cấp định hướng quan trọng cho thiết kế
tính năng trực quan hóa trong eDorm. Thay vì phục vụ mục đích lựa chọn và đặt
phòng trực tiếp, tính năng này được sử dụng để giúp sinh viên hình dung rõ hơn
về môi trường lưu trú trước khi nộp hồ sơ đăng ký và tham gia quá trình xét
duyệt.


\subsubsection*{Khảo sát nền tảng University Living}

University Living (\textit{universityliving.com}) là nền tảng hỗ trợ tìm kiếm chỗ ở dành cho sinh viên quốc tế,
được xây dựng theo hướng cá nhân hóa quá trình lựa chọn nơi lưu trú. Thay vì
hiển thị danh sách phòng ngay từ đầu, hệ thống dẫn dắt người dùng qua quy trình
gồm năm bước liên tiếp: lựa chọn trường học, thời gian nhập học, khoảng cách
mong muốn đến trường, loại phòng và danh sách gợi ý phù hợp. Thông qua việc
thu thập các thông tin này, nền tảng có thể sàng lọc và đề xuất những lựa chọn
đáp ứng tốt hơn nhu cầu của từng sinh viên trước khi hiển thị kết quả.

Mỗi cơ sở lưu trú được giới thiệu bằng nhiều hình thức nội dung khác nhau như
hình ảnh thực tế, video, đánh giá từ sinh viên đã trải nghiệm và tham quan không
gian trực tuyến 360 độ. Cách trình bày này giúp người dùng có cái nhìn toàn diện
hơn so với việc chỉ quan sát ảnh chụp thông thường. Bên cạnh đó, hệ thống còn
hỗ trợ tính toán khoảng cách từ nơi ở đến trường theo từng phương thức di chuyển
như đi bộ, xe đạp hoặc phương tiện công cộng, giúp sinh viên đánh giá tính thuận
tiện của vị trí một cách thực tế hơn.

Ngoài các thông tin mô tả, nền tảng còn khai thác hiệu quả yếu tố cộng đồng
thông qua việc hiển thị số lượng sinh viên đã quan tâm đến một cơ sở lưu trú
hoặc các nhãn đánh giá nổi bật dành cho những lựa chọn được nhiều sinh viên yêu
thích. Những thông tin này góp phần tăng độ tin cậy và hỗ trợ người dùng đưa ra
quyết định phù hợp hơn.

\begin{figure}[H]
\centering
\includegraphics[width=0.9\linewidth]{Hinhve/khaosat_uniliving.png}
\caption{Giao diện khảo sát từ nền tảng University Living}
\label{fig:uniliving}
\end{figure}

Điểm mạnh của University Living nằm ở khả năng cá nhân hóa theo ngữ cảnh: thay
vì để người dùng tự lọc trong một danh sách dài, hệ thống chủ động thu thập
thông tin và trình bày kết quả phù hợp. Cách tiếp cận này làm giảm đáng kể
nhận thức về sự phức tạp của quy trình, ngay cả với người dùng lần đầu sử dụng
hoặc không thành thạo ngôn ngữ địa phương. Đây là bài học quan trọng về thiết
kế quy trình đăng ký: thay vì đổ tải nhận thức lên người dùng, hệ thống nên
dẫn dắt từng bước rõ ràng.

Quy trình thu thập thông tin ban đầu gồm năm bước của University Living hoạt
động hiệu quả trong bối cảnh người dùng chưa xác định rõ nhu cầu và cần hệ
thống hỗ trợ thu hẹp phạm vi lựa chọn. Tuy nhiên, bài toán đăng ký ký túc xá
tại HUST có đặc thù khác biệt khi sinh viên không tìm kiếm và lựa chọn phòng
theo sở thích cá nhân mà tham gia vào một đợt xét duyệt dựa trên các tiêu chí
ưu tiên đã được quy định trước. Vì vậy, trải nghiệm người dùng của eDorm không
tập trung vào việc dẫn dắt khám phá như University Living mà ưu tiên cung cấp
thông tin minh bạch về điều kiện xét duyệt, trạng thái hồ sơ và kết quả phân
bổ, giúp sinh viên dễ dàng theo dõi tiến trình đăng ký và hiểu được cơ sở của
các quyết định phân phòng.

\subsubsection*{Phân tích xu hướng thị trường}

Từ kết quả khảo sát hai nền tảng trên kết hợp với tổng quan các hệ thống quản
lý chỗ ở sinh viên hiện đại, có thể nhận diện bốn xu hướng phát triển nổi bật.

Trực quan hóa không gian lưu trú đang trở thành một thành phần quan trọng trong
trải nghiệm người dùng. Thay vì chỉ cung cấp thông tin dưới dạng văn bản, nhiều
hệ thống hiện đại tích hợp ảnh toàn cảnh 360 độ, sơ đồ tầng và mô hình không
gian nhằm giúp sinh viên hình dung môi trường sinh hoạt trước khi đăng ký.

Số hóa toàn diện quy trình đăng ký và quản lý hồ sơ đang trở thành chuẩn tối
thiểu. Các hệ thống tiên tiến không dừng ở mức số hóa biểu mẫu giấy mà xây dựng
toàn bộ quy trình từ khai báo, nộp minh chứng, xét duyệt đến thông báo kết quả
trong một nền tảng thống nhất, kèm theo cơ chế lưu vết nhằm bảo đảm tính minh
bạch và khả năng giải trình.

Ứng dụng di động và thiết kế đa thiết bị đã chuyển từ tùy chọn sang yêu cầu bắt
buộc. Sinh viên kỳ vọng có thể truy cập thông tin và thực hiện các thao tác cơ
bản trên điện thoại thông minh với trải nghiệm tương đương máy tính.

Tự động hóa phân bổ kết hợp khả năng giải thích kết quả đang nhận được nhiều sự
quan tâm. Thách thức quan trọng không chỉ nằm ở thuật toán mà còn ở khả năng
trình bày kết quả theo cách người dùng có thể kiểm chứng và hiểu được. Minh bạch
trong quy trình phân bổ vì vậy trở thành yếu tố quyết định mức độ tin cậy của
hệ thống.

Thông báo thời gian thực đang dần thay thế mô hình cập nhật theo từng đợt hoặc
theo lịch cố định. Việc cung cấp thông tin ngay khi hồ sơ thay đổi trạng thái
hoặc khi kết quả được công bố giúp sinh viên chủ động theo dõi quá trình đăng
ký, đồng thời giảm đáng kể khối lượng công việc hỗ trợ từ phía cán bộ quản lý.


\subsubsection*{Kết luận thị trường}

Nhìn chung, các nền tảng quốc tế đã giải quyết tốt bài toán tự phục vụ và minh
bạch thông tin, với giao diện hiện đại, đa thiết bị và quy trình dẫn dắt người
dùng rõ ràng. Tuy nhiên, toàn bộ các nền tảng này phục vụ mô hình thị trường
tự do --- sinh viên tự lựa chọn chỗ ở theo khả năng tài chính --- trong khi ký
túc xá đại học công lập tại Việt Nam vận hành theo cơ chế phân bổ tập trung
có điều kiện, với các tiêu chí ưu tiên về học lực, khu vực địa lý, hoàn cảnh
kinh tế và đối tượng chính sách. Đây là sự khác biệt căn bản khiến không thể
áp dụng trực tiếp các giải pháp thương mại trên mà cần một hướng thiết kế riêng,
vừa kế thừa các điểm mạnh về trải nghiệm người dùng, vừa đáp ứng đặc thù vận
hành của môi trường giáo dục công lập Việt Nam. Nhận ra khoảng trống này là
lúc tôi xác định được hướng đi cho eDorm --- không cần cạnh tranh với Scape về
độ bóng bẩy của giao diện, không cần rộng hơn University Living về kho dữ liệu
phòng trọ toàn cầu, mà cần giải đúng bài toán mà cả hai đều không đặt ra: phân
bổ phúc lợi minh bạch, có thể kiểm chứng, và thích ứng được với chính sách thay
đổi mỗi năm.

\subsection{So sánh các giải pháp}

{\hbadness=10000
\begin{table}[htbp]
\centering
\caption{So sánh các giải pháp trong bài toán quản lý phân bổ ký túc xá}
\label{tab:comparison}
\renewcommand{\arraystretch}{1.5}
\begin{tabular}{p{4.0cm} p{4.8cm} p{4.8cm}}
\toprule
\textbf{Giải pháp} & \textbf{Ưu điểm} & \textbf{Nhược điểm} \\
\midrule
Bảng tính điện tử
  & Triển khai nhanh, chi phí thấp, phù hợp thử nghiệm nhỏ.
  & Khó mở rộng, dễ sai sót, thiếu khả năng audit và mô phỏng. \\
Phần mềm quản lý KTX thương mại
  & Ổn định, nhiều tính năng quản trị, hỗ trợ báo cáo.
  & Chi phí cao, hạn chế tùy biến chính sách, tích hợp phức tạp. \\
Giải pháp nghiên cứu (matching/ILP)
  & Cơ sở lý thuyết vững, tối ưu mục tiêu cụ thể.
  & Khó chuyển sang sản phẩm, lời giải khó diễn giải cho người dùng. \\
Hệ thống tích hợp (Engine + UI)
  & Tùy biến chính sách, mô phỏng, realtime, trải nghiệm tốt.
  & Cần đầu tư phát triển, yêu cầu tích hợp dữ liệu từ nhiều nguồn. \\
\bottomrule
\end{tabular}
\end{table}
}% end \hbadness group

Bảng~\ref{tab:comparison} tổng hợp sự đánh đổi giữa các hướng tiếp cận và làm
rõ lý do tại sao không có giải pháp đơn lẻ nào đáp ứng đồng thời toàn bộ yêu
cầu về chất lượng thuật toán, khả năng vận hành và trải nghiệm người dùng. Do
đó, giải pháp trong đề tài này là kết hợp có chọn lọc: một engine với nền tảng
lý thuyết đủ vững để đảm bảo chất lượng phân bổ, tích hợp trong một hệ thống
có đầy đủ cơ chế audit, mô phỏng và giao diện thân thiện.

\subsection{Định hướng cho eDorm}

Từ kết quả khảo sát thực tế và phân tích xu hướng thị trường, có thể suy ra hai
nhóm yêu cầu tính năng đối với hệ thống eDorm: nhóm tính năng bắt buộc phản ánh
những gì đã trở thành chuẩn tối thiểu trên thị trường, và nhóm tính năng tạo
khác biệt thể hiện những điểm chưa được triển khai tốt trong bối cảnh đại học
Việt Nam.

Nhóm tính năng bắt buộc bao gồm: cổng đăng ký trực tuyến để sinh viên tự khai
báo, nộp minh chứng và theo dõi trạng thái hồ sơ mà không cần đến trực tiếp;
thông báo thời gian thực về kết quả phân bổ và thay đổi trạng thái hồ sơ; tra
cứu công nợ và lịch sử thanh toán; bảng điều khiển quản trị tổng hợp tình trạng
phòng ở cho cán bộ; giao diện responsive tương thích đa thiết bị; và ứng dụng
di động cho sinh viên.

Nhóm tính năng tạo khác biệt tập trung vào những điểm mà các nền tảng quốc tế
đã triển khai hiệu quả nhưng hầu như chưa có mặt trong các hệ thống ký túc xá
tại Việt Nam. Trực quan hóa không gian phòng thông qua mô hình 3D và sơ đồ tầng
tương tác giúp sinh viên đưa ra quyết định đăng ký có căn cứ thay vì chỉ dựa
vào thông tin văn bản. Minh bạch hóa điểm ưu tiên --- hiển thị chi tiết điểm
theo từng tiêu chí kèm lời giải thích có thể kiểm chứng --- giải quyết trực tiếp
nguyên nhân chính gây khiếu nại sau mỗi đợt phân bổ. Công cụ mô phỏng phân bổ
cho cán bộ quản lý, cho phép đánh giá tác động của các kịch bản chính sách trước
khi công bố kết quả chính thức, là tính năng gần như chưa từng xuất hiện trong
các hệ thống ký túc xá đại học công lập trong nước. Ngoài ra, nhật ký hoạt động
(audit trail) cấu trúc ghi lại toàn bộ thao tác kèm thời gian thực hiện đáp ứng
yêu cầu giải trình ngày càng cao trong quản lý giáo dục công lập.

Từ những yêu cầu được xác định ở đây, một số quyết định kiến trúc đã phát sinh
khá tự nhiên: cần một engine phân bổ tách biệt khỏi giao diện để có thể chạy
độc lập trong chế độ mô phỏng; cần cơ chế cấu hình chính sách động để cán bộ
tự thay đổi mà không đụng vào code; và cần một lớp thông báo đủ tin cậy để
sinh viên không bao giờ bỏ lỡ kết quả quan trọng. Những quyết định này và lý
do đằng sau chúng là chủ đề của chương tiếp theo.


% =====================================================================
\section{Tổng quan chức năng}
\label{sec:functional}
% =====================================================================

Hệ thống phục vụ bốn nhóm tác nhân chính: sinh viên (đăng ký, nộp minh chứng,
xem kết quả phân bổ, tham quan phòng VR, yêu cầu bảo trì), quản lý KTX (cấu
hình chính sách, vận hành phiên phân bổ, duyệt minh chứng, quản lý sinh viên),
hệ thống tự động (phân bổ tự động, ghi audit log, gửi thông báo realtime) và
nhóm kiểm duyệt (tra cứu báo cáo và thống kê).

Quy trình xử lý của hệ thống gồm các bước: thu thập và chuẩn hoá hồ sơ sinh
viên, cấu hình chính sách phân bổ, chạy mô phỏng để đánh giá các kịch bản
trước khi commit kết quả chính thức, sau đó gửi thông báo đến sinh viên và hỗ
trợ xử lý khiếu nại; việc tách biệt giữa mô phỏng và commit giúp kiểm soát
chất lượng và hạn chế sai sót trong quá trình phân phòng.

\subsection{Biểu đồ use case tổng quát}

Biểu đồ use case tổng quát cung cấp cái nhìn hệ thống về các chức năng cốt lõi
và mối quan hệ giữa các tác nhân chính. Biểu đồ này có vai trò xác lập ranh
giới nghiệp vụ của hệ thống, đồng thời chỉ ra những nhóm chức năng mức cao sẽ
được phân rã ở các mục tiếp theo.

\begin{figure}[H]
\centering
\includegraphics[width=0.98\linewidth]{Hinhve/usecase_overview.png}
\caption{Biểu đồ use case tổng quát của hệ thống quản lý ký túc xá}
\label{fig:usecase_overview}
\end{figure}

Biểu đồ tổng quát (Hình~\ref{fig:usecase_overview}) xác định ba tác nhân chính
gồm sinh viên, quản lý KTX và hệ thống tự động. Sinh viên thực hiện các hành
động liên quan tới đăng ký, nộp minh chứng, xem xếp hạng và kết quả phân bổ,
yêu cầu bảo trì, đổi/trả phòng, tham quan phòng VR, nhận thông báo và xem hồ
sơ cá nhân. Quản lý KTX chịu trách nhiệm quản lý đơn đăng ký, duyệt minh chứng
ưu tiên, quản lý kỳ đăng ký, cấu hình chính sách phân bổ, vận hành phiên mô
phỏng và thực thi phân bổ, quản lý KTX và phòng, quản lý sinh viên, xử lý yêu
cầu bảo trì, xem báo cáo thống kê và gửi thông báo. Hệ thống tự động đảm nhiệm
phân bổ tự động, ghi audit log và gửi thông báo realtime --- trong đó chức năng
tham quan phòng VR và phân bổ tự động đều bao gồm (include) việc xem mô hình 3D
phòng ở trước khi đăng ký.

\subsection{Biểu đồ use case phân rã}
Nhóm biểu đồ phân rã được trình bày tập trung trong một sơ đồ duy nhất nhằm
minh hoạ mối quan hệ phân tầng giữa các chức năng mức cao và các chức năng con.
Sơ đồ này thể hiện tám nhóm chức năng mức cao cùng các thành phần phân rã tương
ứng.
\begin{figure}[H]
\centering
\includegraphics[width=0.92\textwidth]{Hinhve/eDorm-WBS.png}
\caption{Biểu đồ phân rã chức năng (toàn cảnh)}
\label{fig:phanra}
\end{figure}
Phân tích phân rã chức năng (tóm tắt):

Xác thực \& Tài khoản phân rã thành đăng nhập đa phương thức (MSSV, Google,
HUST Portal), đăng ký tài khoản, xác minh OTP qua email và bảo mật 2FA kèm
quản lý hồ sơ; nhóm này đảm bảo định danh người dùng đúng đắn trước khi truy
cập các chức năng nghiệp vụ.

Tra cứu \& Hỗ trợ phân rã thành tra cứu KTX và phòng, tham quan VR Tour
$360^\circ$, thẻ cư trú QR và gửi yêu cầu bảo trì; nhóm này tối ưu cho tính
minh bạch thông tin và tương tác trực tiếp của sinh viên với hệ thống.

Đăng ký KTX phân rã thành nộp đơn đăng ký, khai báo hồ sơ ưu tiên, theo dõi
kết quả và kiểm duyệt hồ sơ bởi quản trị viên; các chức năng con này đảm bảo
tính chuẩn hoá dữ liệu đầu vào và tạo nền tảng cho bước xét ưu tiên tiếp theo.

Cấu hình \& Chính sách phân rã thành quản lý kỳ đăng ký, thiết lập chính sách
và trọng số ưu tiên, quản lý quota và kiểm soát phiên bản kèm kiểm toán; đây
là vùng chịu trách nhiệm chính về cách thức tính điểm ưu tiên và ràng buộc
phân bổ.

Mô phỏng \& Phân bổ phân rã thành chạy mô phỏng, so sánh phương án, xác nhận
commit và phân tích cohort kèm dự báo; chức năng này tách rõ hai chế độ preview
và commit, đồng thời hỗ trợ ra quyết định chính sách dựa trên dữ liệu lịch sử.

Quản lý phòng phân rã thành các chức năng thêm/sửa/xoá phòng, cập nhật sức
chứa, liên kết ảnh và VR Tour, và xét duyệt chuyển phòng; đảm bảo dữ liệu
không gian luôn chuẩn xác để engine phân bổ sử dụng.

Thông báo \& Realtime phân rã thành đẩy thông báo, hàng đợi retry, ghi nhật
ký sự kiện và cấu hình kênh gửi; đảm bảo tính kịp thời và tin cậy trong giao
tiếp với người dùng.

Báo cáo \& Thống kê phân rã thành báo cáo phân bổ, thống kê KPI, dự báo nhu
cầu và xuất dữ liệu định dạng CSV/XLSX; nhóm này hỗ trợ phân tích kết quả vận
hành và ra quyết định chính sách dài hạn.

\subsection{Quy trình nghiệp vụ}

Quy trình nghiệp vụ tổng thể của hệ thống được thiết kế theo chuỗi xử lý liên
tục, bắt đầu từ tiếp nhận hồ sơ đăng ký và kết thúc ở giai đoạn thông báo,
phản hồi và xử lý các phát sinh sau phân bổ. Trước hết, sinh viên thực hiện
đăng ký và nộp minh chứng; sau đó hệ thống kiểm tra dữ liệu đầu vào và chuyển
hồ sơ hợp lệ sang trạng thái sẵn sàng xét duyệt. Trên cơ sở dữ liệu đã chuẩn
hóa, cán bộ quản lý cấu hình chính sách và tham số phân bổ, từ đó khởi chạy
mô phỏng để đánh giá tác động của cấu hình hiện hành. Khi kết quả mô phỏng phù
hợp, hệ thống được chuyển sang chế độ thực thi để ghi nhận kết quả chính thức.

Sau khi kết quả phân bổ được công bố, sinh viên tra cứu kết quả, đọc phần giải
thích và phản hồi nếu phát sinh khiếu nại. Đồng thời, hệ thống phát thông báo
thời gian thực tới các bên liên quan và ghi lại audit log cho toàn bộ thao tác
quan trọng.

% =====================================================================
\section{Đặc tả chức năng}
\label{sec:usecase}
% =====================================================================

Phần này trình bày năm use case trọng yếu đại diện cho các luồng nghiệp vụ
chính của hệ thống. Mỗi use case được mô tả theo khuôn mẫu chuẩn bao gồm tác
nhân, điều kiện tiên quyết, điều kiện hậu kết và luồng thực hiện chính. Hình
minh hoạ tương ứng cho mỗi use case được chèn ngay sau phần mô tả.

\subsection{Use Case 1: Đăng ký chỗ ở và nộp minh chứng ưu tiên}

Sinh viên có nhu cầu ở ký túc xá đăng nhập vào hệ thống, điền thông tin cá
nhân và lựa chọn KTX/loại phòng mong muốn. Song song đó, sinh viên tải lên
các tài liệu minh chứng ưu tiên nếu có (PDF, JPG, PNG, tối đa 10~MB mỗi tệp).
Hệ thống kiểm tra tính hợp lệ của dữ liệu và tệp đính kèm: nếu sai định dạng
hoặc thiếu thông tin bắt buộc, thông báo lỗi được trả về và yêu cầu tải lại.
Khi hồ sơ hợp lệ, hệ thống tính điểm ưu tiên sơ bộ, lưu hồ sơ vào cơ sở dữ
liệu và gửi xác nhận đến sinh viên.

\textbf{Tác nhân:} Sinh viên (chính), Hệ thống (phụ).

\textbf{Tiền điều kiện:} Sinh viên đã có tài khoản hợp lệ được cấp bởi nhà trường và đang trong thời hạn mở đăng ký.

\textbf{Hậu điều kiện:} Hồ sơ đăng ký ở trạng thái ``đã tiếp nhận'', điểm ưu tiên sơ bộ được tính, hồ sơ sẵn sàng cho bước xét duyệt minh chứng và phân bổ.

\textbf{Luồng phát sinh:} Tệp không đúng định dạng hoặc vượt giới hạn dung lượng --- hệ thống từ chối và yêu cầu tải lại; thông tin bắt buộc còn thiếu --- hệ thống thông báo cụ thể và dừng quy trình.

\begin{figure}[H]
\centering
\includegraphics[width=0.99\textwidth]{Hinhve/usecase_registration.png}
\caption{Use Case 1 --- Đăng ký chỗ ở và nộp minh chứng ưu tiên}
\label{fig:uc_registration}
\end{figure}

\subsection{Use Case 2: Cấu hình chính sách và quản lý kỳ đăng ký}

Quản trị viên thực hiện hai nhóm công việc liên quan đến việc chuẩn bị cho một
đợt phân bổ. Nhóm đầu tiên là quản lý kỳ đăng ký: mở hoặc đóng đợt đăng ký,
thiết lập thời hạn nộp hồ sơ, và duyệt minh chứng ưu tiên từng hồ sơ (xem chi
tiết minh chứng, phê duyệt hoặc từ chối kèm lý do). Minh chứng được phê duyệt
sẽ có hiệu lực trong phiên phân bổ tiếp theo. Nhóm thứ hai là cấu hình chính
sách phân bổ: chọn/tạo mẫu chính sách, thiết lập trọng số tiêu chí ưu tiên,
thiết lập ràng buộc sức chứa và phân nhóm. Mọi cấu hình đều được kiểm tra hợp
lệ tham số trước khi lưu (chặn nếu tổng trọng số $\neq$ 100\%, sức chứa bằng
0 hoặc âm, thiếu ràng buộc giới tính). Phiên bản hoá cấu hình cho phép so sánh
và khôi phục về cấu hình trước.

\textbf{Tác nhân:} Quản trị viên (chính), Hệ thống (phụ).

\textbf{Tiền điều kiện:} Tài khoản có vai trò quản lý KTX với đầy đủ quyền cấu hình.

\textbf{Hậu điều kiện:} Cấu hình chính sách và kỳ đăng ký được lưu vào cơ sở dữ liệu, sẵn sàng kích hoạt cho phiên phân bổ; minh chứng được duyệt có hiệu lực điểm ưu tiên.

\textbf{Luồng phát sinh:} Tham số vi phạm ràng buộc logic --- hệ thống báo lỗi chi tiết và yêu cầu điều chỉnh trước khi lưu.

\begin{figure}[H]
\centering
\includegraphics[width=0.99\textwidth]{Hinhve/usecase_config.png}
\caption{Use Case 2 --- Cấu hình chính sách và quản lý kỳ đăng ký}
\label{fig:uc_config}
\end{figure}

\subsection{Use Case 3: Mô phỏng và thực thi phân bổ}

Đây là use case trung tâm của hệ thống, với điều kiện tiên quyết quan trọng là
quản trị viên đã chạy ít nhất một phiên mô phỏng và xác nhận kịch bản trước khi
commit.

Ở \textbf{chế độ Preview}, quản trị viên khởi chạy Allocation Engine trên tập
dữ liệu thực nhưng không ghi kết quả vào cơ sở dữ liệu chính thức. Hệ thống
sinh báo cáo tổng quan về kết quả phân bổ dự kiến: tỷ lệ lấp đầy phòng, phân
bổ điểm ưu tiên theo nhóm, danh sách các trường hợp không thể phân bổ do thiếu
phòng hoặc vi phạm ràng buộc. Quản trị viên có thể so sánh nhiều kịch bản với
tham số khác nhau trước khi quyết định.

Ở \textbf{chế độ Commit}, quản trị viên xác nhận kịch bản đã chọn. Hệ thống
khởi chạy lại Allocation Engine lần cuối, cập nhật trạng thái phòng và hồ sơ
sinh viên, ghi audit log đầy đủ và gửi thông báo kết quả tới sinh viên. Sinh
viên sau đó có thể xem kết quả phân bổ cá nhân.

\textbf{Tác nhân:} Quản trị viên (chính), Hệ thống (phụ), Sinh viên (nhận thông báo kết quả).

\textbf{Tiền điều kiện:} Dữ liệu ứng viên và cấu hình chính sách đã sẵn sàng; đã chạy ít nhất một phiên mô phỏng và admin xác nhận kịch bản trước khi commit.

\textbf{Hậu điều kiện:} Kết quả phân bổ chính thức được lưu trữ, có thể tra cứu và xuất báo cáo; thông báo đã gửi đến sinh viên.

\textbf{Luồng phát sinh:} Engine phát hiện vi phạm ràng buộc cứng không thể giải quyết --- hệ thống liệt kê chi tiết các trường hợp và đề xuất điều chỉnh tham số.

\begin{figure}[H]
\centering
\includegraphics[width=0.98\textwidth]{Hinhve/usecase_simulation.png}
\caption{Use Case 3 --- Mô phỏng và thực thi phân bổ}
\label{fig:uc_simulation}
\end{figure}

\subsection{Use Case 4: Quản lý KTX và danh sách phòng}

Quản trị viên thực hiện quản lý ở hai cấp độ. Ở \textbf{cấp tòa nhà (KTX)},
các thao tác bao gồm thêm/sửa/vô hiệu hoá KTX, cập nhật thông tin và sức chứa
tổng, liên kết ảnh và mô hình VR định dạng \texttt{.glb} (được nén Draco xuống
dưới 20~MB để phục vụ VR tour trên cả web và mobile). Từ màn hình quản lý KTX,
quản trị viên có thể xem tình trạng sinh viên cư trú hiện tại. Ở \textbf{cấp
phòng đơn lẻ}, các thao tác bao gồm thêm/sửa/vô hiệu hoá phòng, cập nhật sức
chứa và loại phòng, gán trang thiết bị và tiện nghi. Trước khi lưu, hệ thống
kiểm tra hợp lệ dữ liệu và từ chối nếu sức chứa nhỏ hơn 1, thiếu tên hoặc loại
phòng, hoặc trùng số phòng trong cùng tầng.

\textbf{Tác nhân:} Quản trị viên (chính), Hệ thống (phụ).

\textbf{Tiền điều kiện:} Tài khoản có vai trò quản trị hệ thống.

\textbf{Hậu điều kiện:} Danh sách KTX và phòng được cập nhật nhất quán; mô hình VR được liên kết và sẵn sàng cho sinh viên tham quan; dữ liệu sức chứa chuẩn để engine phân bổ sử dụng.

\textbf{Luồng phát sinh:} Dữ liệu vi phạm ràng buộc (sức chứa âm, trùng số phòng, thiếu loại phòng) --- hệ thống từ chối và trả về thông báo lỗi cụ thể.

\begin{figure}[H]
\centering
\includegraphics[width=0.99\textwidth]{Hinhve/usecase_rooms.png}
\caption{Use Case 4 --- Quản lý KTX và danh sách phòng}
\label{fig:uc_rooms}
\end{figure}

\subsection{Use Case 5: Thông báo thời gian thực và audit log}

Hệ thống tự động phát hiện các sự kiện quan trọng: kết quả phân bổ được công
bố, trạng thái hồ sơ thay đổi, có cập nhật liên quan đến yêu cầu bảo trì, hoặc
sinh viên gửi đánh giá dịch vụ. Khi đánh giá từ 2 sao trở xuống, cảnh báo ưu
tiên cao được gửi đến quản trị viên. Với mỗi sự kiện được phát hiện, hệ thống
thực hiện đồng thời ba hành động: gửi thông báo in-app qua WebSocket (chỉ áp
dụng cho người dùng đã cài app và cấp quyền thông báo), gửi push notification
qua Expo API đến thiết bị di động, và ghi audit log không thể sửa/xoá vào bộ
lưu trữ riêng biệt. Sinh viên có thể xem danh sách thông báo trong ứng dụng;
quản trị viên có thể truy cập nhật ký kiểm toán đầy đủ để tra cứu mọi sự kiện
hệ thống.

\textbf{Tác nhân:} Hệ thống (chính), Sinh viên và Quản trị viên (nhận kết quả).

\textbf{Tiền điều kiện:} Dịch vụ WebSocket hoạt động bình thường; Expo Push Notification service khả dụng.

\textbf{Hậu điều kiện:} Thông báo đã được gửi đến thiết bị của người dùng qua ít nhất một kênh; audit log được ghi vào bộ lưu trữ riêng với thuộc tính immutable (không thể sửa/xoá).

\textbf{Luồng phát sinh:} WebSocket bị gián đoạn --- hệ thống chuyển sang hàng đợi retry; người dùng chưa cấp quyền thông báo --- bỏ qua push notification, chỉ ghi log.

\begin{figure}[H]
\centering
\includegraphics[width=0.99\textwidth]{Hinhve/usecase_notifications.png}
\caption{Use Case 5 --- Thông báo thời gian thực và audit log}
\label{fig:uc_notifications}
\end{figure}

% =====================================================================
\section{Yêu cầu phi chức năng}
\label{sec:nonfunctional}
% =====================================================================

Các yêu cầu phi chức năng dưới đây không phải là danh sách lý tưởng được chọn
từ tài liệu kỹ thuật --- chúng xuất phát từ bối cảnh cụ thể của hệ thống. Mỗi
đợt đăng ký KTX tập trung hàng trăm sinh viên truy cập đồng thời trong khung
giờ ngắn, nên hiệu năng dưới tải cao không phải là tùy chọn. Dữ liệu sinh viên
bao gồm hoàn cảnh gia đình và minh chứng ưu tiên là thông tin nhạy cảm, nên bảo
mật phải được tích hợp từ đầu, không phải vá sau. Cán bộ vận hành không phải
là nhân viên IT, nên hệ thống phải đủ ổn định để không cần can thiệp kỹ thuật
trong các đợt phân bổ thực tế.

Về hiệu năng, hệ thống phải xử lý được tối thiểu 5.000 yêu cầu đọc đồng thời,
bao gồm tra cứu trạng thái và hiển thị kết quả, với độ trễ P95 dưới 300 ms
trong điều kiện tải bình thường. Đối với các tác vụ ghi phức tạp như chạy phiên
phân bổ hay cập nhật đồng loạt, hệ thống hỗ trợ ít nhất 200 yêu cầu đồng thời.
Các tác vụ nặng như chạy Allocation Engine trên tập dữ liệu lớn được thực hiện
theo mô hình bất đồng bộ với thông báo khi hoàn thành, tránh block giao diện
người dùng.

Về bảo mật, toàn bộ API được bảo vệ bằng JWT kết hợp RBAC với ba vai trò chính:
cán bộ quản lý toàn quyền, cán bộ chỉ xem và sinh viên. Thông tin cá nhân và
tài liệu minh chứng được mã hoá khi lưu trữ. Mọi giao tiếp giữa client và
server sử dụng TLS/HTTPS. Audit log cho các thao tác quan trọng được lưu trong
bộ lưu trữ riêng với quyền ghi một lần, đảm bảo tính bất biến.

Về độ tin cậy và sẵn sàng, mục tiêu availability là 99,9\% cho các chức năng
cốt lõi trong giờ vận hành bình thường. Dữ liệu được sao lưu tự động theo lịch
và quy trình khôi phục được kiểm thử định kỳ. Hệ thống được thiết kế để khởi
động lại gracefully mà không mất dữ liệu giao dịch đang xử lý.

Về khả năng mở rộng, kiến trúc được thiết kế để hỗ trợ mở rộng theo chiều ngang
cho tầng API khi lưu lượng tăng. Cơ sở dữ liệu MongoDB Atlas hỗ trợ sharding
khi kích thước collection vượt ngưỡng hiệu năng. Các thành phần được đóng gói
bằng Docker để triển khai nhất quán trên mọi môi trường.

Về tính dễ sử dụng và truy cập, giao diện sinh viên tuân theo nguyên tắc
mobile-first, đảm bảo trải nghiệm tốt trên màn hình nhỏ với băng thông giới
hạn. Giao diện hỗ trợ đầy đủ các trình duyệt phổ biến và không yêu cầu cài đặt
plugin hay phần mềm bổ sung.

Về tính dễ bảo trì, Allocation Engine và các service lõi phải đạt coverage kiểm
thử đơn vị tối thiểu 80\%. Log hệ thống có cấu trúc JSON để dễ tích hợp với các
công cụ phân tích log. Mã nguồn tuân theo các quy ước đặt tên và tổ chức thư
mục nhất quán, kèm tài liệu API theo chuẩn OpenAPI 3.0.

Về tương thích và stack công nghệ, các lựa chọn bao gồm Node.js/Express cho
backend triển khai dạng Docker container; MongoDB Atlas cho cơ sở dữ liệu đám
mây; Redis cho tầng realtime adapter; Socket.IO cho thông báo thời gian thực;
EJS kết hợp Bootstrap cho giao diện web; Expo/React Native cho ứng dụng di
động. Thiết kế collection cơ sở dữ liệu bao gồm trên 19 schema với chỉ mục phù
hợp cho tra cứu theo mã sinh viên, trạng thái và thời gian.

Những yêu cầu trên không phải là danh sách nguyện vọng --- chúng là cam kết có
thể đo được. Một số trong đó, đặc biệt hiệu năng dưới tải 300 người dùng đồng
thời, đã được kiểm chứng bằng load test thực tế trong quá trình phát triển. Sự
khác biệt giữa ``yêu cầu đặt ra'' và ``yêu cầu đã kiểm chứng'' là điều quan
trọng trong vận hành thực tế --- và đó là một trong những lý do chương kiểm thử
(chương 4) dành riêng một mục cho kết quả load test.

Từ những vấn đề quan sát được trong chương này --- sự thiếu minh bạch, thiếu
công cụ mô phỏng, và khoảng cách giữa bài toán phúc lợi với các giải pháp thị
trường hiện có --- điều tôi tin chắc là hướng thiết kế được chọn đúng: không
phải xây thêm một cổng đăng ký nữa, mà xây một hệ thống mà cả sinh viên lẫn
cán bộ đều thực sự hiểu được điều gì đang xảy ra bên trong nó. Chương 3 trình
bày những lựa chọn công nghệ cụ thể để hiện thực hoá hướng đi đó.

\end{document}